\section{El panorama China–Estados Unidos: cómo superar el bloqueo geopolítico}

La sección anterior trató el mapa de empresas; esta cierra con el panorama entre China y Estados Unidos. Al final de la entrevista se insiste en que la inversión y el emprendimiento tecnológicos no pueden basarse en moverse dentro de los círculos de contactos, sino en crear. El bloqueo geopolítico limita el cómputo, el capital, los mercados y la cooperación, pero no determina por sí solo quién gana; lo que realmente atraviesa esas restricciones es la creatividad técnica, la ejecución de ingeniería, la comprensión del producto, la difusión mediante código abierto y la capacidad de organizar la cadena industrial.

\begin{figure}[H]
\centering
\includegraphics[width=0.92\textwidth]{china-us-geofence.png}
\caption{El panorama de la IA entre China y Estados Unidos: bajo el bloqueo geopolítico, la creatividad técnica importa más que moverse en los círculos de contactos. Diagrama conceptual propio, elaborado a partir de la conversación de 01:54:32--02:01:11.}
\end{figure}

\begin{knowledgebox}{Lectura de la figura: tener ventajas distintas no obliga a limitarse a la defensa}
Del lado estadounidense están el cómputo, los frontier lab, el capital y los puntos de acceso a productos globales; del lado chino, la velocidad de ingeniería, los escenarios de aplicación, la eficiencia del código abierto y la capacidad de ejecución de la cadena industrial. El bloqueo geopolítico genera fricción, pero también obliga a los equipos a buscar nuevos caminos en eficiencia, ingeniería y ecosistema local.
\end{knowledgebox}

\subsection{Dónde están las oportunidades de los equipos chinos}

Esta sección divide las oportunidades en cuatro tipos. La primera es de eficiencia: construir modelos y productos competitivos con menos cómputo. La segunda es de aplicación: China tiene una gran cantidad de escenarios de negocio digitalizado, contenido, comercio electrónico, industria y hardware. La tercera es de código abierto: obtener distribución entre desarrolladores de todo el mundo mediante modelos y herramientas abiertos. La cuarta es de cadena industrial: cuando la IA se combina con hardware, robots, automóviles y electrónica de consumo, la cadena de suministro y la velocidad de ingeniería de China cobran importancia.

\begin{importantbox}{Experiencia práctica: en tiempos de bloqueo, hace más falta una creación verificable}
Cuanto más fuertes son las restricciones externas, menos basta con el discurso. Los equipos necesitan presentar modelos, código, productos, ingresos, curvas de eficiencia y adopción por parte de desarrolladores que puedan verificarse. La verdadera influencia global proviene de lo que se crea, no del respaldo de un círculo.
\end{importantbox}

\subsection{Riesgos bajo las restricciones geopolíticas}

Esta sección también debe mantener la conciencia de los riesgos. Los equipos chinos enfrentan limitaciones en el suministro de cómputo, el cumplimiento normativo en el extranjero, la confianza en la marca, las ventas a empresas, los pagos y el acceso a los ecosistemas; los equipos estadounidenses también enfrentan costos, regulación, burbujas de producto, expectativas del capital y una creciente complejidad organizacional. El panorama China–Estados Unidos no es una presión unilateral ni un rebase unilateral, sino dos sistemas de innovación que buscan avances bajo restricciones distintas.

\noindent\begin{tabularx}{\textwidth}{p{0.18\textwidth}p{0.34\textwidth}X}
\toprule
Dimensión & Ventajas/riesgos de EE.\,UU. & Ventajas/riesgos de China \\
\midrule
Cómputo & Suministro y ecosistema de nube sólidos, pero con alta presión de costos y regulación & Limitado por los controles de exportación; requiere eficiencia y alternativas. \\
Modelos & Alta concentración de frontier lab & La eficiencia del código abierto y las recipe de ingeniería pueden ser la vía de avance. \\
Productos & Puntos de acceso globales y clientes empresariales sólidos & Escenarios locales abundantes, pero una marca global es más difícil. \\
Cadena industrial & Ecosistema de software sólido & Cadenas sólidas de hardware, automóviles, robots y manufactura. \\
\bottomrule
\end{tabularx}

\subsection{Resumen de la sección}

Lo esencial del panorama de la IA entre China y Estados Unidos no es un simple optimismo o pesimismo, sino ver quién logra crear mejores modelos, productos, entornos y ecosistemas bajo restricciones. Para seguir a largo plazo el ámbito de la IA e internet, conviene fijarse en productos verificables y no solo en el entusiasmo que genera el discurso.
